%% =====================================================================
%%  COVER LETTER
%% =====================================================================
%%  The letterhead, date, recipient block, subject line, salutation,
%%  sign-off and enclosure list are generated from config.tex.
%%  You only write the body.
%% =====================================================================
\begin{coverletter}[
  % date       = {November 12, 2026},  % default: today's date
  % closing    = {Sincerely,},
  % subject    = none,                 % `auto' (default), `none' or your own text
  % enclosures = auto,                 % `auto' lists the other documents; `none' or text
  % page limit = 2,                    % warns you if the letter runs longer
]

\tip[How committees read this letter]{Aim for 1.5 to 2 pages for
research-intensive posts. Some fields expect a single page, so check the
norms of yours. A reliable arc is: who you are and what you are applying
for; your record; your research vision and how you will fund it; why this
department; teaching and mentoring; a short close. Write as a future
colleague, not as a supplicant. Do not repeat your CV; interpret it.}

I am delighted to apply for the position of \PositionTitle in
\PositionField in the \Department at \Institution. I am an
environmental data scientist who builds trustworthy forecasting tools for
freshwater ecosystems by combining process-based models, machine learning
and community-collected data. I am currently an NSERC Postdoctoral Fellow at
the Northbridge Institute of Technology. I completed my PhD in Environmental
Science at the University of Eastvale in 2024. My work has produced nine
peer-reviewed publications (six as first author) and an open-source
early-warning system that two provincial monitoring agencies now use.

\textbf{Research record.} Harmful algal blooms threaten drinking water,
recreation and fisheries, yet most lakes are monitored too sparsely to see a
bloom coming. My doctoral research showed that hybrid models can forecast bloom
onset ten days ahead with 35\% lower error than either process-based or
purely data-driven models. These hybrids embed the physics of lake mixing
inside a neural network. The work was published in the \emph{Journal of
Freshwater Informatics} and received the journal's Editor's Choice
designation. As a postdoctoral fellow, I turned these methods into
\emph{BloomCast}, an open-source system that issues probabilistic forecasts
for 42 lakes each day. Agency staff use it to schedule sampling, which
cut the median time to detect a bloom from nine days to three. Building
BloomCast taught me to design models around the decisions people actually
make. That lesson now drives my research agenda.

\tip[Show, don't tell]{Replace adjectives with evidence: numbers, venues,
users, funding, outcomes. At research-intensive institutions research
usually fills 40--50\% of the letter. Teaching-focused institutions
weight teaching much more heavily, so adjust the balance.}

\textbf{Research vision.} My lab will ask how to make environmental forecasts
reliable when the future no longer resembles the past. Over the next five
years I will pursue three connected aims:
\begin{enumerate}[label=(\arabic*)]
  \item \emph{forecasting under change}: methods that detect and adapt to
        climate-driven shifts in lake behaviour;
  \item \emph{community data with community control}: integrating
        volunteer and Indigenous-led monitoring into forecasts under data
        agreements that keep ownership with the communities that collect
        the data;
  \item \emph{decision-ready forecasts}: co-designing forecast products with
        water managers and testing whether they improve real decisions.
\end{enumerate}
I will fund this program through an NSERC Discovery Grant in my first
year, an NSERC Alliance partnership with provincial agencies and
conservation authorities, and a CFI John R.\ Evans Leaders Fund application
for a real-time sensor network. I have already contributed to a successful
\$480{,}000 Alliance grant and hold a \$25{,}000 seed grant as principal
investigator. My research statement sets out this program in detail.

\tip[Tailor this paragraph to each job]{This paragraph is the one that
shows you chose this department. Name specific colleagues, centres,
facilities, courses and strategic priorities. Keep everything that changes
between applications here and in \texttt{config.tex}, and adapting the
letter takes minutes.}

\textbf{Why \Institution.} The \Department is an exceptional fit. The
Westbrook Centre for Freshwater Futures and its 30-year monitoring record
for the Westbrook River watershed would let my lab test forecasts
across a whole watershed rather than lake by lake. I look forward to
collaborating with Dr.~Priya Raman on coupling watershed hydrology with lake
forecasts, and with Dr.~Tom\'as Ortega on the ecological consequences of
blooms. I would also be glad to help build the \Faculty's new minor in
Environmental Data Science, which matches my teaching experience closely.

\textbf{Teaching and mentorship.} I have taught or co-taught two
undergraduate courses and supported four more as a teaching assistant,
reaching more than 450 students. I redesigned \emph{Introduction to
Environmental Data Analysis} around real datasets and universal design for
learning. In the redesigned course, the share of students who withdrew or
failed fell by half, and students rated my teaching 4.6 out of 5. At
\Institution I could teach quantitative methods, hydrology and data science
courses immediately. I would also develop an upper-year course in machine
learning for environmental systems and a graduate course in environmental
forecasting. I have mentored fourteen student researchers, five of whom are
co-authors on my publications, and I co-founded a paid summer research
program for students from groups under-represented in the environmental
sciences. My teaching dossier and EDIA statement describe this work in
detail.

Thank you for considering my application. I would welcome the chance to
discuss how my research and teaching could contribute to the \Department.

\end{coverletter}
